\hypertarget{chapter-33-c23-std-print}{%
\section{CHAPTER 33: C23 STD PRINT}\label{chapter-33-c23-std-print}}

\hypertarget{c23-the-end-of-iostream-and-printf}{%
\section{C++23: THE END OF IOSTREAM AND
PRINTF}\label{c23-the-end-of-iostream-and-printf}}

\hypertarget{stdprint-and-stdprintln}{%
\subsubsection{\texorpdfstring{1. \texttt{std::print} and
\texttt{std::println}}{1. std::print and std::println}}\label{stdprint-and-stdprintln}}

C++23 finally fixes standard output. \passthrough{\lstinline!std::cout!}
is slow and verbose, while \passthrough{\lstinline!printf!} is not
type-safe. \passthrough{\lstinline!std::print!} bridges the gap using
the C++20 \passthrough{\lstinline!std::format!} engine.

\begin{itemize}
\item
  \textbf{Type-safe and Fast}: It writes directly to the underlying OS
  file descriptor without creating an intermediate
  \passthrough{\lstinline!std::string!} allocation.

\begin{lstlisting}[language={C++}]
#include <print>

std::println("User {} has {} points.", user.name, user.score);
std::println(stderr, "Error: connection failed");
\end{lstlisting}
\end{itemize}

\hypertarget{formatting-ranges}{%
\subsubsection{2. Formatting Ranges}\label{formatting-ranges}}

\passthrough{\lstinline!std::print!} natively understands standard
ranges and containers.

\begin{lstlisting}[language={C++}]
std::vector<int> v = {1, 2, 3};
std::println("{}", v); // Output: [1, 2, 3]
\end{lstlisting}
